% Chapter 4 -- Organizing Your Files (Part I, Learning; AN-14 task form).
% Sources (Law 1): os/milton/kmain.c flair_live_do_surface (a click
% selects, Shift-click adds to the selection, a click on bare surface
% deselects, a press-drag on bare surface draws the selection rectangle);
% os/flair/finder_windows.c (fw_do_new_folder: NEWFOLD, then NEWFOL02 to
% NEWFOL99 in the front window's folder; fw_do_select_all; Clean Up snaps
% to the 68 x 52 grid; fw_do_arrange_by_name: case-insensitive name order;
% the window keeps each icon's place); os/flair/finder_ops.c (drop targets:
% a folder icon, the INITECH icon, another disk window; refusals zoom the
% outline back: service, samedir, stale, cycle, exists, locked);
% os/flair/finder_windows.h (FINDER_WIN_MAX 4 windows, 64 icons a window,
% the hidden TRASH and DESKTOP.DB).  Model: Macintosh User's Guide 1991
% Ch. 6 "Organizing Your Work" (apple-macintosh.md Sec 8.2).
% GATES: test-flair-disk-windows test-flair-file-ops test-flair-desktop-icons test-finder-windows test-finder-ops test-flair-finder-cmds test-pubs-flair-screens
\Gates{test-flair-disk-windows, test-flair-file-ops, test-flair-desktop-icons,
  test-finder-windows, test-finder-ops, test-flair-finder-cmds,
  test-pubs-flair-screens}
\chapter{Organizing Your Files}
\label{ch:organizing}

The information you keep on a disk is stored in \term{files}: documents
you create, and the applications that create them.  You keep related files
together in \term{folders}, and folders can hold other folders.  In FLAIR
each file and folder is shown as an icon, and you organize your files by
moving their icons.

\section{Looking at the files in a folder}
\index{folders!opening}
\begin{TaskSteps}
\Step Open the \texttt{INITECH} icon, if its window is not already open.
\Step Double-click the \file{APPS} folder in the \texttt{Drive A Files}
window.
\Result{The \texttt{APPS} window opens.  It holds the applications and the
programs that come with FLAIR.}
\end{TaskSteps}
\Operation{FINDER-OPEN-FOLDER}

\FlairScreen[label=fig:appswin]{appswin}{The APPS folder}

The icons tell you what kind of thing each one is.

\SideHead{Application}A diamond: a program you can open from the
desktop, such as \file{TEXT.EXE}, InitechText.  Its name ends in
\file{.EXE}.

\SideHead{Document}A page with a folded corner: a file of information,
such as \file{README.TXT}, or a program that runs only at the InitechDOS
prompt, such as \file{CHKDSK.COM}.

\SideHead{Folder}A folder: a container for other files and folders.

Names are shown as InitechDOS stores them: up to eight characters, a
period and up to three more.  A name too long for the space under its icon
is cut short on the screen, as \file{TENANTFX.EXE} is in
Figure~\ref{fig:appswin}.

\section{Selecting more than one icon}
\index{selecting!several icons}
Commands and drags act on everything that is selected.  There are three
ways to select more than one icon in a window.

\SideHead{Shift-click}Click the first icon.  Then hold down \key{Shift}
and click each of the others.

\FlairScreen[label=fig:shiftclick]{shiftclick}{Two icons selected with the
Shift key}

\SideHead{Selection rectangle}Point at a bare part of the window, hold
down the mouse button and drag diagonally.  A rectangle follows the
pointer; when you release the button, every icon inside the rectangle is
selected.

\FlairScreen[label=fig:marquee]{marquee}{Icons selected by dragging a
rectangle round them}

\SideHead{Select All}Choose \menuitem{Select All} from the
\menuname{Edit} menu, or press \key{Ctrl-A}.  Every icon in the active
window is selected; with no Finder window open, every icon on the desktop
is selected.

\FlairScreen[label=fig:selectall]{selectall}{Select All in the Drive A
Files window}
\Operation{FINDER-SELECT-ALL}
\Operation{FINDER-WIN-MARQUEE}
\Operation{FINDER-WIN-SELECT}

To deselect everything, click a bare part of the window.

\section{Making a new folder}
\index{folders!new}\index{New Folder command}
\Operation{FINDER-NEW-FOLDER}
\begin{TaskSteps}
\Step Click in the \texttt{Drive A Files} window to make it active.
\Result{The new folder is made in the folder of the active window.}
\Step Choose \menuitem{New Folder} from the \menuname{File} menu.
\Result{A folder named \file{NEWFOLD} appears in the window, after the
icons already there.}
\end{TaskSteps}

\FlairScreen[label=fig:newfolder]{newfolder}{A new folder}

If a folder named \file{NEWFOLD} already exists, the new folder is named
\file{NEWFOL02}, then \file{NEWFOL03}, and so on up to \file{NEWFOL99}.
The name of a folder cannot be changed in FLAIR.

\section{Moving an icon into a folder}
\index{moving files}\index{folders!moving files into}
\Operation{FINDER-MOVE}
\begin{TaskSteps}
\Step Drag the \file{README.TXT} icon onto the \file{NEWFOLD} folder.
\Result{When the pointer is over the folder, the folder is darkened to
show that it will receive what you release.}
\Step Release the button.
\Result{\file{README.TXT} disappears from the \texttt{Drive A Files}
window: the file has moved into the folder.  Its place in the window
stays empty until you tidy the window or open it again.}
\Step Double-click \file{NEWFOLD}.
\Result{Its window opens, with \file{README.TXT} inside.}
\end{TaskSteps}

\FlairScreen[label=fig:intofolder]{intofolder}{README.TXT moved into the new
folder}

You can move an icon in the same way onto a folder icon in any open
window, into the body of another open Finder window, or onto the
\texttt{INITECH} icon, which moves it to the top level of the data disk.
Moving a file changes only where it is kept; the file itself is not
changed, and no copy is made.

When FLAIR cannot make a move, the outline of the icon slides back to where
it started and nothing changes.  That happens when:
\begin{itemize}
\item the place you release it already holds a file or folder of the same
  name;
\item you drag a folder into itself or into a folder inside it;
\item you release the icon in the folder it is already in, or on nothing
  that can receive it.
\end{itemize}
A locked file can be moved from folder to folder like any other, but it
cannot be put in the Trash (Chapter~\ref{ch:trash}).

\section{Moving an icon within a window}
\index{icons!moving}
\Operation{FINDER-WIN-DRAG}
Drag an icon to a bare part of its own window to move it there.  The icon
stays where you put it until the window is closed.  Each time a Finder
window opens, its icons are laid out in order in rows from the top left;
only the icons on the desktop itself keep the places you give them.

\FlairScreen[label=fig:strayicon]{strayicon}{README.TXT moved within the
window}

\section{Tidying a window}
\index{Clean Up command}\index{Arrange command}
Two commands put the icons of the active window back in order.

\SideHead{Clean Up}Choose \menuitem{Clean Up} from the \menuname{View} or
the \menuname{Special} menu.  Each icon moves to a place of the window's
invisible grid, so that the icons stand in tidy rows.
\Operation{FINDER-CLEANUP}

\FlairScreen[label=fig:cleanedup]{cleanedup}{The window after Clean Up}

\SideHead{Arrange}Choose \menuitem{Arrange (by Name)} from the
\menuname{View} menu.  The icons are laid out row by row in alphabetical
order of their names, filling any gaps.
\Operation{FINDER-ARRANGE}

\FlairScreen[label=fig:arranged]{arranged}{The window after Arrange (by
Name)}

\section{Limits}
\index{limits}
FLAIR shows at most 64 icons in one window, and opens at most four Finder
windows at a time.  If a window does not open, close one of the others
first.  The Trash folder and the file \file{DESKTOP.DB}, in which FLAIR
records the arrangement of your desktop, are on the data disk but are not
shown in its window.
